\begin{annexitem}{annex:p3.morley.codes}{thm:p03-morley-code-space}{A complete metric on the interpretation coordinates}
\begin{proof}
Each relation or function contributes countably many
coordinates, since its arity is finite; a constant
contributes one. The language is countable, so all
coordinates together can be listed in a finite or
countable sequence. Each alphabet is a nonempty finite
or countable discrete space. Choose one default value
in each coordinate. The codes differing from defaults
in only finitely many coordinates form a countable
dense set.

For coordinates numbered by \(j<\omega\), the metric
\[
 d(u,v)=\sum_{j<\omega}2^{-j-1}
                         \mathbf1_{u(j)\ne v(j)}
\]
induces the product topology. Unused coordinates in a
finite list can be given singleton alphabets.
A Cauchy sequence eventually has a fixed value in
each coordinate, because changing coordinate \(j\)
costs at least \(2^{-j-1}\).
The coordinatewise limit is a valid code.
Agreement on a sufficiently long initial segment makes
the tail sum arbitrarily small, proving convergence.
Thus the metric is complete.
The same construction applies to a finite domain;
an empty coordinate list gives a singleton space.

On the countable topological sum, use the corresponding
bounded complete metric within each component and
distance \(2\) between different components.
Every Cauchy sequence eventually remains in one component,
so this metric is complete. The union of the component
dense sets is countable and dense.

The space of countable strings in a countable alphabet
has cardinality at most \(|\omega^\omega|=2^{\aleph_0}\).
For completeness, binary strings inject into it, and
each \(\omega\)-sequence has a graph that is a subset
of the countable set \(\omega^2\), giving the reverse
upper bound. A countable sum of such code spaces still
has this upper bound.
Every nonempty at most countable structure has a labeled
copy in one component. Selecting representatives or
passing to the quotient cannot increase this bound.
The fixed-finite-domain example in the perfect-family
proof below supplies the asserted possible continuum
of finite isomorphism types.
\end{proof}
\end{annexitem}

\begin{annexitem}{annex:p3.morley.borel-truth}{thm:p03-morley-borel-truth}{Term evaluation and the full well-founded syntax induction}
\begin{proof}
For a fixed assignment, evaluation of each finitary
term is a continuous map from codes to the discrete
domain. A variable is fixed; a constant reads one
coordinate. For \(f(t_1,\ldots,t_k)\), induction
makes each argument value locally fixed near any code.
Reading the function at this fixed argument tuple
then makes its result locally fixed.
This proves continuity, even when nested functions
read different coordinates at different codes.
An atomic equality or relation truth value is likewise
locally fixed. Both its truth and falsity sets are open,
so the atomic truth set is clopen.

Induct on the well-founded syntax tree of a fixed
infinitary formula. Negation takes a complement;
countable conjunction takes a countable intersection;
countable disjunction takes a countable union.
For existential quantification, the truth set at
the fixed old assignment is
\[
 \bigcup_{b\in\omega}
       \{M:M\models\psi(\bar a,b)\}.
\]
Universal quantification gives the corresponding
countable intersection.
All these operations preserve Borelness.
The finite-domain case uses the same induction,
with finite quantifier unions and intersections.

For assignments varying in \(\omega^n\), partition
the product into its countably many clopen assignment
slices. Satisfaction is the union of the fixed-assignment
Borel truth sets on these slices, so it is Borel.
For \(n=0\) there is one slice.

Each formula has countably many syntax-tree nodes.
Indeed every node has at most countably many children
and finite-distance levels are countable; the full
tree is their countable union.
Well-founded induction on it and regularity of
\(\omega_1\) give a countable quantifier rank under
the displayed recursion.
A countable conjunction of sentences has no free
variables and a well-founded countable syntax tree;
it is a permitted sentence. For an empty family use
\(\top\), and for finite families use the ordinary
finite conjunction. Satisfaction of this sentence is
exactly satisfaction of the entire family, giving
the asserted Borel model class.
No uncountable family is conjoined by this argument.
\end{proof}
\end{annexitem}

\begin{annexitem}{annex:p3.morley.approximants}{thm:p03-morley-approximant-properties}{Transitivity, the Borel recursion and formula preservation}
\begin{proof}
At stage \(0\), agreement on atomic truth values is
an equivalence relation. Suppose the stage-\(\alpha\)
relations are equivalence relations in every arity.
Reflexivity at \(\alpha+1\) chooses the same new point;
symmetry is built into the two extension clauses.
For transitivity, let pointed \(M,N,P\) have the
two successive stage-\(\alpha+1\) correspondences.
A new point \(c\) of \(M\) has a matching \(d\)
in \(N\), and that \(d\) has a matching \(e\)
in \(P\). The induction hypothesis on their
enlarged tuples gives the matching stage-\(\alpha\)
correspondence from \(M\) to \(P\).
The reverse extension starts in \(P\) and passes
through \(N\) to \(M\).
The old tuples are stage-\(\alpha\)-equivalent
by transitivity there.
Thus they are stage-\(\alpha+1\)-equivalent.
At a limit, intersection preserves the equivalence
relation properties. Each successor retains the
preceding stage, and limits retain all earlier ones;
transfinite induction proves decrease with the ordinal.

For fixed coded tuples at stage \(0\), matching the
truth value of each atomic formula is a clopen
condition on the pair of structure codes.
There are countably many atomic formulas because
there are countably many symbols and finitary terms.
Their conjunction is closed, hence Borel; it need
not be clopen for an infinite atom list.
For the successor, a forth extension condition is
a countable intersection over source points of
countable unions over target points of the preceding
relation on enlarged tuples. The back condition is
the reverse one.
Their intersection with the old relation is Borel.
A countable limit ordinal has only countably many
predecessors, so its intersection is also Borel.
The tuple coordinates range over countable discrete
sets; their slice union supplies the uniform pointed
relation. This proves every fixed-ordinal Borel claim.

Preservation of formulas is a simultaneous
well-founded syntax induction. If an approximant
level is at least an atom's rank, it implies the
stage-\(0\) match. Negation transfers the induction
assertion in both directions. For a conjunction or
disjunction, its rank is at least the rank of each
constituent; decreasing the approximant to those
levels transfers every constituent and hence the
connective. This includes countable connectives.
For \(\exists y\,\psi\), use the approximant at
\(\operatorname{qr}(\psi)+1\).
A source witness has a target match at
\(\operatorname{qr}(\psi)\), so the induction
hypothesis transfers \(\psi\).
The back clause gives the reverse implication.
Renaming the bound variable to a fresh slot handles
an existential originally binding one of the old
variable symbols. The new coordinate is appended
to the tuple; unused old slots do not affect truth.
Universal quantification is the dual case.
This proves preservation at every ordinal at least
the formula's rank.
\end{proof}
\end{annexitem}

\begin{annexitem}{annex:p03-morley-scott-stabilization}{thm:p03-morley-scott-stabilization}{Countably many possible first separations}
\begin{proof}
The atomic conjunction is countable and decides
exactly the source tuple's atomic truth values.
Thus its target satisfaction is exactly the
stage-\(0\) approximant.
Inductively, the old formula in a successor retains
the old tuple match. Its conjunction of existential
formulas supplies a target witness for every
source extension. Its universally quantified
disjunction supplies a source match for every
target extension. By the induction hypothesis
these are exactly the two stage-\(\alpha\)
extension clauses. At a limit the conjunction
gives all preceding matches.
This proves the semantic equivalence for every
structure \(N\), not just countable ones.
All formulas remain in the specified logic:
the atomic set, source domain and each countable
ordinal's predecessor set are countable, and the
free-variable set is the fixed finite tuple.

Within \(M\), there are only countably many pairs
of equal-length finite tuples, over all finite
lengths. For each pair that fails some countable
approximant, choose its least failing ordinal.
Let \(\delta_M\) be the supremum of these ordinals,
with supremum \(0\) if there are none.
It is countable by regularity of \(\omega_1\).
Every pair ever separated at a countable stage
is already separated at \(\delta_M\), by decrease.
Therefore a pair agreeing at \(\delta_M\) cannot
first disagree at \(\delta_M+1\): that latter
ordinal would be in the chosen set and at most
its supremum, contradicting agreement at the
supremum. This proves simultaneous stabilization
for every finite arity.
It makes no assertion of a common \(\delta_M\)
as \(M\) varies over uncountably many codes.
\end{proof}
\end{annexitem}

\begin{annexitem}{annex:p3.morley.scott-sentence}{thm:p03-morley-scott-sentence}{The encoded stability clauses and complete onto construction}
\begin{proof}
The family of all finite tuples in \(M\), over all
finite arities, is countable.
Each implication in the display has finitely many
quantifiers and no free variables after its universal
closure. Their conjunction is therefore a valid
infinitary sentence.
The empty-tuple formula holds in \(M\) by reflexivity.
If a tuple \(\bar b\) of \(M\) satisfies the antecedent
attached to \(\bar a\), the semantic equivalence
gives \((M,\bar a)\equiv_{\delta_M}(M,\bar b)\).
Stabilization gives the successor relation, hence
the displayed consequent. Thus \(M\models\sigma_M\).

Suppose an at most countable \(N\) satisfies \(\sigma_M\).
The empty-tuple clause gives an initial
stage-\(\delta_M\) match.
Maintain a finite tuple match
\((M,\bar a)\equiv_{\delta_M}(N,\bar b)\).
By semantic equivalence, \(N\) satisfies the
antecedent attached to this particular source
tuple \(\bar a\). Its universal stability clause
then gives the successor formula there.
The existential clauses provide a
stage-\(\delta_M\) match for each new source point;
the universal-disjunction clause provides a source
match for each new target point.
This is the required extension in both directions.
The stage-\(0\) equality match identifies exactly
the same repeated coordinates.
If a new source point is outside the old source
entry set, its inequalities to all old entries
force the target match outside the old target
set. The reverse statement holds for a new
target point. The tuple correspondence therefore
extends a finite injective partial map.

Enumerate the at most countable source and target
domains. Alternate adding the first not-yet-covered
source point and the first not-yet-covered target
point, skipping a direction once its finite domain
has been exhausted.
The extension clauses prevent an unmatched point
on one side when every point on the other has
already been matched: such a new point would
require a fresh mate by the inequalities.
The finite or infinite union is consequently
a bijection between the two domains.
Every finite collection of arguments appears in
one stage, where all atomic formulas agree.
Relations are therefore preserved and reflected.
For a function, include its argument tuple and
its value together in one stage; the atomic
equation for that value proves preservation of
the function. The same reasoning with \(x=c\)
preserves constants.
Thus the union is an isomorphism.
An isomorphic copy satisfies \(\sigma_M\) because
isomorphisms preserve infinitary formulas by
the same syntax induction, proving the converse.

After naming a fixed finite tuple, the expanded
language is still countable and the construction
applies unchanged. An isomorphism must preserve
each named constant, so it has the specified
tuple image. Repeated names have equal
interpretations on both sides.
Countability was used for the onto enumeration
and the countable tuple conjunction; the result
does not identify arbitrary uncountable models.
\end{proof}
\end{annexitem}

\begin{annexitem}{annex:p3.morley.approximant-isomorphism}{thm:p03-morley-approximant-isomorphism}{Every Scott sentence has a countable formula rank}
\begin{proof}
An isomorphism matches atomic values and all
one-point extensions in both directions.
Induction gives each approximant, including
the empty-tuple relations.
Conversely suppose all empty-tuple approximants
hold. The Scott sentence \(\sigma_M\) has a
countable quantifier rank by the proved syntax
rank result. The approximant at that rank
preserves its truth.
As \(M\models\sigma_M\), also \(N\models\sigma_M\),
and the full Scott construction makes \(M,N\)
isomorphic.
The Borel recursion supplied each individual
\(E_\alpha\). There are \(\aleph_1\) countable
ordinals, not a countable enumeration of all
of them. Borel closure under countable
intersection cannot be applied to this
full intersection.
\end{proof}
\end{annexitem}

\begin{annexitem}{annex:p3.morley.isomorphism-complexity}{thm:p03-morley-isomorphism-complexity}{Scott truth for one class and a closed permutation graph}
\begin{proof}
For a fixed \(M\), its isomorphism class on
domain \(\omega\) is exactly the truth set
of \(\sigma_M\). It is Borel by the
infinitary-truth lemma.

To verify the Polish witness space directly,
give permutations the topology of pointwise
convergence. Inversion is continuous: requiring
\(\pi^{-1}(i)=j\) is exactly requiring
\(\pi(j)=i\). A compatible complete metric is
\[
 d(\pi,\tau)=\sum_{i<\omega}2^{-i-2}
  \left(\mathbf1_{\pi(i)\ne\tau(i)}
       +\mathbf1_{\pi^{-1}(i)\ne\tau^{-1}(i)}\right).
\]
A Cauchy sequence stabilizes every forward
and inverse coordinate. Their limits are
mutual inverses: if the forward limit takes
\(i\) to \(j\), eventual coordinate agreement
forces the inverse limit to take \(j\) to \(i\).
Thus the limit is a permutation, and the
tail estimate proves metric convergence.
Finite-support permutations form a countable
dense subset. Indeed a finite consistent
partial bijection extends to a permutation
of the finite union of its domain and range,
and fixes the remaining points.
This also matches any finite collection of
forward and inverse coordinate conditions.
Hence \(S_\infty\) is Polish.

Let \(C\) consist of triples \((\pi,M,N)\)
where \(\pi\) is an isomorphism from \(M\) to \(N\).
For each relation and each source argument tuple,
require its source truth value to equal the
target truth value on the image tuple.
For each function and argument tuple require
\[
 \pi(f^M(\bar a))=f^N(\pi\bar a),
\]
and for each constant require
\(\pi(c^M)=c^N\).
Every individual requirement is clopen:
near any triple only finitely many discrete
values need to be fixed to determine the
two sides. Dependence on evaluated input
coordinates is locally fixed, exactly as
in the term-evaluation proof.
There are countably many requirements.
Their intersection \(C\) is closed.
Its projection to \((M,N)\) is precisely
isomorphism, since the first coordinate
already ranges over all bijections.
This proves analyticity under the stated
projection definition.
The argument does not make the projection
closed or Borel; nor does one Borel truth
set for each \(M\) form a countable uniform
list of all classes.
\end{proof}
\end{annexitem}

\begin{annexitem}{annex:p3.morley.perfect-example}{thm:p03-morley-perfect-example}{Every branch, closedness and separation by predicates}
\begin{proof}
Every code \(M_\eta\) satisfies \(U\), because
each predicate is uniformly true or false.
If \(\eta\ne\nu\), choose \(i\) where they
differ. One model has \(P_i\) everywhere and
the other nowhere. Since domains are nonempty,
no isomorphism can preserve this predicate
between them. The family is pairwise
nonisomorphic.

The interpretation coordinate for \(P_i(a)\)
is exactly \(\eta(i)\), so fixing finitely
many code coordinates imposes only finitely
many binary-coordinate conditions.
This proves continuity of the branch map.
Its inverse on the image reads \(P_i(0)\)
for each \(i\), and is continuous as well.
The image is exactly the set of codes
satisfying
\[
 P_i(a)=P_i(b)\quad\text{for every }i,a,b<\omega,
\]
where equality here compares Boolean
interpretation values. Each condition is
clopen and their intersection is closed.
The image is nonempty.
In any basic neighborhood only finitely
many predicate bits are specified; changing
a different bit gives a distinct point
in that neighborhood. Hence there are
no isolated points. The image is perfect.

There are \(2^{\aleph_0}\) branches and all
give distinct isomorphism types. The code
upper bound gives the reverse inequality,
proving the exact spectrum.
The sentences \(\forall xP_i(x)\) have
different truth values in the family, so
\(U\) is not complete.
On a fixed finite nonempty domain the same
all-or-none interpretations remain
pairwise nonisomorphic for the identical
predicate reason.
Thus even fixed finite domains have
continuum many types of structures in
this countable infinite language.
Neither family purports to realize many
models of one complete first-order theory.
\end{proof}
\end{annexitem}
